\documentclass[14pt]{extarticle}

% ===== Русский язык и шрифты (pdfLaTeX) =====
\usepackage[T2A]{fontenc}
\usepackage[utf8]{inputenc}
\usepackage[russian]{babel}
\addto\captionsrussian{\renewcommand{\contentsname}{Оглавление}}
\usepackage{mathptmx}
\usepackage{caption}
\usepackage{array}
\usepackage{tabularx}
\usepackage{longtable}
\usepackage{xcolor}
\usepackage{float}
\usepackage{geometry}
\usepackage{setspace}
\usepackage{indentfirst}
\usepackage{graphicx}
\usepackage{tocloft}
\usepackage{fvextra}
\usepackage{titlesec}
\usepackage{fancyhdr}
\usepackage{tikz}
\usetikzlibrary{arrows.meta,positioning,shapes.geometric}
\usepackage{microtype}
\microtypesetup{expansion=false}
\usepackage{xurl}
\emergencystretch=3em

% ===== Колонтитулы =====
\pagestyle{fancy}
\fancyhf{}
\fancyfoot[R]{\thepage}
\renewcommand{\headrulewidth}{0pt}
\fancypagestyle{plain}{%
  \fancyhf{}%
  \fancyfoot[R]{\thepage}%
  \renewcommand{\headrulewidth}{0pt}%
}

% ===== Поля, интервалы, абзацы =====
\geometry{
  a4paper,
  left=3cm,
  right=1.5cm,
  top=2cm,
  bottom=2cm
}
\onehalfspacing
\setlength{\parindent}{0.8cm}

% ===== Картинки =====
\graphicspath{{images/}}

\captionsetup[figure]{
    name={Рисунок},
    labelsep={space},
    justification=centering,
    singlelinecheck=false
}

% ===== Оглавление =====
\renewcommand{\cftsecleader}{\cftdotfill{\cftdotsep}}
\renewcommand{\cftsecfont}{\normalsize}
\renewcommand{\cftsecpagefont}{\normalsize}
\renewcommand{\cftsubsecfont}{\normalsize}
\renewcommand{\cftsubsecpagefont}{\normalsize}
\renewcommand{\cftsubsubsecfont}{\normalsize}
\renewcommand{\cftsubsubsecpagefont}{\normalsize}
\renewcommand{\cfttoctitlefont}{\hfill\large\bfseries}
\renewcommand{\cftaftertoctitle}{\hfill}
\setlength{\cftbeforesecskip}{4pt}
\setlength{\cftbeforesubsecskip}{2pt}
\setlength{\cftbeforesubsubsecskip}{1pt}

% ===== Код =====
\DefineVerbatimEnvironment{CodeBlock}{Verbatim}{
  breaklines=true,
  breakanywhere=true,
  fontsize=\footnotesize,
  baselinestretch=1,
  formatcom=\color{black}
}

% ===== Таблицы =====
\newcolumntype{Y}{>{\raggedright\arraybackslash}X}
\renewcommand{\arraystretch}{1.2}

% ===== Заголовки =====
\titleformat{\section}
  {\normalfont\normalsize\bfseries}
  {\thesection}{1em}{}

\titleformat{\subsection}
  {\normalfont\normalsize\bfseries}
  {\thesubsection}{1em}{}

\titleformat{\subsubsection}
  {\normalfont\normalsize\bfseries}
  {\thesubsubsection}{1em}{}

\titlespacing*{\section}{0.8cm}{12pt}{1.5cm}
\titlespacing*{\subsection}{0.8cm}{12pt}{12pt}
\titlespacing*{\subsubsection}{0.8cm}{12pt}{6pt}

\begin{document}

\begin{titlepage}
\thispagestyle{empty}
\begin{center}
Министерство транспорта Российской Федерации\\
Федеральное агентство железнодорожного транспорта\\[0.2em]

Федеральное государственное бюджетное образовательное учреждение\\
высшего образования\\
«Дальневосточный государственный университет путей сообщения»\\[0.2em]

Кафедра «Цифровая индустрия и защита информации»
\vfill

{\bfseries
\large ОТЧЁТ О ТЕХНОЛОГИЧЕСКОЙ\\[0.3em]
\large (НАУЧНО-ИССЛЕДОВАТЕЛЬСКОЙ) ПРАКТИКЕ
}\\[0.3em]
\Large НИР 09.03.02. 01.01.БО231ИСТ
\large
\vfill

\begin{center}
\begin{minipage}{\textwidth}
  \setlength{\tabcolsep}{0pt}
  \begin{tabular}{@{}p{4cm}p{9cm}p{4cm}@{}}
    Студент
      & \centering\hrulefill
      & \centering Т.П.~Чигирёв \\[0em]
  \end{tabular}
\end{minipage}
\end{center}

\vspace{-25pt}
\small\textit{(подпись, дата)}
\large

\begin{center}
\begin{minipage}{\textwidth}
  \setlength{\tabcolsep}{0pt}
  \begin{tabular}{@{}p{4cm}p{9cm}p{4cm}@{}}
    Руководитель
      & \centering\hrulefill
      & \centering Е.В.~Сазанова \\[0em]
  \end{tabular}
\end{minipage}
\end{center}

\vspace{-20pt}
\small\textit{(подпись, дата)}
\normalsize

\vfill
\normalsize Хабаровск 2026

\end{center}
\end{titlepage}

\setcounter{page}{2}

\tableofcontents
\clearpage

\section*{Введение}
\addcontentsline{toc}{section}{Введение}

Во время научно-исследовательской практики выполнялась работа с учебным полигоном, который имитировал инфраструктуру предприятия с несколькими сетевыми сегментами, веб-сервисами, сервером удалённого доступа, системой мониторинга Cacti и сервером баз данных PostgreSQL. Практическая часть заключалась в анализе инцидента информационной безопасности, выявлении использованных векторов атаки, устранении последствий и описании порядка реагирования для администратора информационной безопасности.

В ходе выполнения задания были последовательно разобраны три связанные уязвимости: компрометация SSH-сервера, эксплуатация узла Cacti через тестовое устройство и компрометация PostgreSQL вследствие небезопасных параметров сетевого доступа и аутентификации. Работа выполнялась по шагам: подключение к узлам полигона, проверка признаков закрепления злоумышленника, анализ журналов PostgreSQL, восстановление удалённой базы данных из снимка, изменение конфигурационных файлов и регистрация инцидентов в системе Ampire.

В данном отчёте систематизированы выполненные действия и оформлена последовательность мер, которая может использоваться при повторении аналогичной атаки. Пароли из инфраструктурного листа в отчёте не приводятся; вместо этого указаны факты смены учётных данных, корректировки правил доступа и восстановления сервисов.

\clearpage

\section*{1. Сведения о предприятии прохождения практики}
\addcontentsline{toc}{section}{1. Сведения о предприятии прохождения практики}

% TODO: раздел заполняется студентом самостоятельно.

\clearpage

\section*{2. Информационные системы предприятия и сетевая инфраструктура}
\addcontentsline{toc}{section}{2. Информационные системы предприятия и сетевая инфраструктура}

\subsection*{2.1 Общая характеристика информационных систем предприятия}
\addcontentsline{toc}{subsection}{2.1 Общая характеристика информационных систем предприятия}

В период практики изучалась типовая инфраструктура предприятия, в которой используются внутренние серверы, сетевые шлюзы, сервисы удалённого администрирования, веб-приложения, система мониторинга и серверы баз данных. Такие информационные системы обеспечивают хранение и обработку рабочих данных, доступ сотрудников к внутренним ресурсам, контроль состояния узлов и регистрацию событий информационной безопасности.

Практическая часть научно-исследовательской работы выполнялась на учебном полигоне. Полигон выступал моделью инфраструктуры предприятия: на нём реальные классы корпоративных систем были заменены виртуальными машинами и сервисами с заранее заданными адресами. Поэтому дальнейшее описание инцидентов относится к учебной среде, но сами меры реагирования соответствуют типовым действиям администратора ИБ на предприятии.

В этой модели сетевая структура была построена по сегментированному принципу: имелись пограничный шлюз, внутренний шлюз, рабочая станция группы реагирования и серверный сегмент. В работе использовались терминал Linux, SSH-клиент, браузер, просмотрщик PDF с перечнем активов, а также штатные средства PostgreSQL. Результаты реагирования фиксировались в системе Ampire.

В практической части были задействованы следующие информационные системы и узлы:

\begin{table}[H]
\caption*{\hfill \textbf{Таблица 1 - Информационные системы предприятия и их представление на полигоне}}
\begin{tabularx}{\textwidth}{|p{4.2cm}|p{3.2cm}|Y|}
\hline
\textbf{Информационная система} & \textbf{Адрес на полигоне} & \textbf{Назначение} \\
\hline
Ampire & 10.11.25.13 & Регистрация карточек инцидентов и фиксация результатов реагирования. \\
\hline
ВМ группы реагирования & 10.11.25.92, 10.11.25.67 & Рабочая станция, с которой открывались удалённые консоли, браузер и инструменты администрирования. \\
\hline
SSH Server & 10.10.1.31 & Узел удалённого доступа; использовался как источник или промежуточная точка атаки. \\
\hline
Cacti & 10.10.2.32 & Система мониторинга; в инциденте фигурирует эксплуатация тестового device и изменение sudoers. \\
\hline
PostgreSQL & 10.10.2.26 & Сервер базы данных; на нём обнаружены следы удалённого доступа, попытки выполнения команд и удаление базы данных. \\
\hline
Edge Gateway & 10.10.1.254 & Пограничный шлюз инфраструктуры. \\
\hline
Internal Gateway & 10.10.1.253 & Внутренний шлюз между сегментами полигона. \\
\hline
\end{tabularx}
\end{table}

\subsection*{2.2 Связь корпоративных ИС с учебным полигоном}
\addcontentsline{toc}{subsection}{2.2 Связь корпоративных ИС с учебным полигоном}

В инфраструктуре предприятия подобные системы выполняют разные функции, но при атаке оказываются связаны между собой. SSH-сервер обеспечивает удалённое администрирование, Cacti используется для мониторинга состояния устройств и сервисов, PostgreSQL хранит данные прикладных систем, а шлюзы разделяют сетевые сегменты. Если один из компонентов настроен небезопасно, злоумышленник может использовать его как промежуточную точку для перехода к другим информационным системам.

На учебном полигоне эта логика была отработана через несколько виртуальных машин. Для расследуемого сценария наибольшее значение имели:
\begin{itemize}
  \item SSH-сервер, так как его компрометация позволяла злоумышленнику получить устойчивый доступ к инфраструктуре и запускать команды от имени локального пользователя;
  \item Cacti, так как система мониторинга взаимодействует с сетевыми устройствами и при небезопасной настройке может стать точкой выполнения команд;
  \item PostgreSQL, так как некорректные настройки сетевого доступа и аутентификации позволяли подключаться к базе данных из внутренней сети и использовать привилегированные функции.
\end{itemize}

\clearpage

\section*{3. Перечень потенциальных векторов атаки}
\addcontentsline{toc}{section}{3. Перечень потенциальных векторов атаки}

\subsection*{3.1 Компрометация SSH-сервера}
\addcontentsline{toc}{subsection}{3.1 Компрометация SSH-сервера}

Первым вектором выделен несанкционированный доступ к узлу SSH Server с адресом 10.10.1.31. Далее этот узел был указан как источник последующих инцидентов в Ampire. При проверке SSH Server выполнялось подключение к серверу, был просмотрен статус пользовательских сервисов и обнаружено задание cron, запускающее файл \texttt{/home/user/evil.conf} каждую минуту.

Вероятной причиной компрометации можно считать использование слабых или известных учётных данных, а также отсутствие достаточных ограничений на сетевой доступ к SSH. Последствием стало закрепление злоумышленника на узле через периодически выполняемое задание и возможность дальнейших действий во внутренней инфраструктуре.

\subsection*{3.2 Эксплуатация Cacti через тестовое устройство}
\addcontentsline{toc}{subsection}{3.2 Эксплуатация Cacti через тестовое устройство}

Второй вектор связан с системой мониторинга Cacti на узле 10.10.2.32. При оформлении инцидента было указано, что злоумышленник использовал тестовый device для получения прав администратора. В качестве индикатора компрометации было зафиксировано изменение файла sudoers.

Для систем мониторинга особенно опасны тестовые или забытые объекты мониторинга. Они часто сохраняют доступ к механизмам опроса, запуску команд, сетевым проверкам и локальным скриптам. Если такой объект доступен злоумышленнику, он может быть использован для выполнения команд на сервере Cacti и повышения привилегий через изменение правил sudo.

\subsection*{3.3 Компрометация PostgreSQL}
\addcontentsline{toc}{subsection}{3.3 Компрометация PostgreSQL}

Третьим вектором выделена атака на узел PostgreSQL с адресом 10.10.2.26. При анализе журналов сервера были обнаружены подключения из внутренней сети, выполнение SQL-команд, обращение к системным объектам и попытка запуска системной команды через механизм \texttt{COPY FROM PROGRAM}. Также было установлено, что была удалена пользовательская база данных \texttt{demo}.

Основными предпосылками атаки стали избыточно открытая сетевая конфигурация PostgreSQL и небезопасное правило аутентификации в \texttt{pg\_hba.conf}. При устранении последствий были отключены широкое правило доступа, изменён параметр \texttt{listen\_addresses} на \texttt{localhost}, обновлён пароль пользователя \texttt{postgres} и восстановлена база данных из снимка.

\clearpage

\section*{4. Термины в области информационной безопасности}
\addcontentsline{toc}{section}{4. Термины в области информационной безопасности}

\begin{longtable}{|p{4.2cm}|p{10.2cm}|}
\caption*{\hfill \textbf{Таблица 2 - Термины, используемые при расследовании инцидента}}\\
\hline
\textbf{Термин} & \textbf{Определение} \\
\hline
Инцидент информационной безопасности & Событие или совокупность событий, которые приводят или могут привести к нарушению конфиденциальности, целостности или доступности информации. \\
\hline
Актив & Сервер, сервис, база данных, учётная запись или иной объект, имеющий ценность для организации и требующий защиты. \\
\hline
Вектор атаки & Способ, которым злоумышленник получает доступ к системе или развивает атаку внутри инфраструктуры. \\
\hline
Индикатор компрометации & Признак, по которому можно установить факт атаки: подозрительная запись в журнале, изменённый файл, сетевое соединение или неизвестный процесс. \\
\hline
SSH & Протокол защищённого удалённого доступа к серверу, который при слабой настройке может стать точкой первоначальной компрометации. \\
\hline
Подбор пароля & Попытка получить доступ к учётной записи путём перебора или использования известных слабых паролей. \\
\hline
Cron & Планировщик задач в Unix-подобных системах; может использоваться злоумышленником для закрепления через регулярный запуск вредоносного файла. \\
\hline
Закрепление в системе & Действия злоумышленника, позволяющие сохранить доступ после перезагрузки или смены сессии. \\
\hline
Повышение привилегий & Получение прав более высокого уровня, например переход от обычного пользователя к администратору. \\
\hline
sudoers & Конфигурация Linux, определяющая, какие пользователи могут выполнять команды через \texttt{sudo}; её изменение является критичным признаком компрометации. \\
\hline
Cacti & Система мониторинга сетевой инфраструктуры, использующая устройства, шаблоны и механизмы опроса для сбора данных. \\
\hline
Device в Cacti & Объект мониторинга, описывающий устройство или хост; при ошибочной настройке может стать каналом выполнения нежелательных действий. \\
\hline
PostgreSQL & Реляционная система управления базами данных, требующая строгого ограничения сетевого доступа и прав пользователей. \\
\hline
\texttt{pg\_hba.conf} & Конфигурационный файл PostgreSQL, определяющий, кто, откуда и каким методом может подключаться к базам данных. \\
\hline
\texttt{postgresql.conf} & Основной конфигурационный файл PostgreSQL, в котором задаются параметры сервера, включая адреса прослушивания. \\
\hline
\texttt{COPY FROM PROGRAM} & Привилегированный механизм PostgreSQL, позволяющий выполнять внешнюю программу при наличии соответствующих прав; при компрометации суперпользователя может быть использован для запуска команд ОС. \\
\hline
Reverse shell & Обратное соединение от атакованного узла к внешнему адресу злоумышленника для получения командной оболочки. \\
\hline
\texttt{pg\_shadow} & Системное представление PostgreSQL с чувствительной информацией об учётных записях; обращение к нему обычным сценарием работы не является типичным. \\
\hline
Снимок базы данных & Резервная копия или SQL-дамп, используемый для восстановления базы после удаления или повреждения данных. \\
\hline
Hardening & Комплекс мер по усилению защищённости системы: отключение лишнего доступа, смена паролей, ограничение прав и проверка журналов. \\
\hline
IDS & Система обнаружения вторжений, предназначенная для выявления подозрительной сетевой или системной активности. \\
\hline
SIEM & Система централизованного сбора, корреляции и анализа событий безопасности. \\
\hline
\end{longtable}

\clearpage

\section*{5. Схема проведения атаки}
\addcontentsline{toc}{section}{5. Схема проведения атаки}

После анализа карточек инцидентов и журналов PostgreSQL сценарий атаки был представлен как последовательная компрометация нескольких узлов. Первоначальная точка была связана с SSH Server, после чего злоумышленник получил возможность использовать внутренние сервисы. Далее атака перешла на Cacti, где тестовое устройство позволило развить атаку и повысить привилегии. После этого был затронут PostgreSQL: сервер принимал подключения из внутренней сети, что позволило выполнить SQL-команды, обратиться к чувствительным системным объектам и удалить базу данных.

\begin{figure}[H]
\centering
\begin{tikzpicture}[
  node distance=1.3cm,
  box/.style={rectangle, draw, rounded corners=2pt, align=center, minimum width=4.5cm, minimum height=1cm, font=\small},
  bad/.style={rectangle, draw, rounded corners=2pt, align=center, minimum width=4.5cm, minimum height=1cm, fill=red!8, font=\small},
  fix/.style={rectangle, draw, rounded corners=2pt, align=center, minimum width=4.5cm, minimum height=1cm, fill=green!8, font=\small},
  arrow/.style={-{Latex[length=2.5mm]}, thick}
]

% Левая колонка - атака
\node[bad] (ssh) at (0,0) {SSH Server 10.10.1.31\\слабый доступ, cron};
\node[bad, below=0.9cm of ssh] (cacti) {Cacti 10.10.2.32\\тестовый device, sudoers};
\node[bad, below=0.9cm of cacti] (pg) {PostgreSQL 10.10.2.26\\открытый доступ, удаление БД};

% Правая колонка - устранение
\node[fix, right=2.5cm of ssh] (sshfix) {Смена пароля\\проверка cron\\ограничение SSH};
\node[fix, right=2.5cm of cacti] (cactifix) {Удаление device\\восстановление sudoers\\перезапуск сервиса};
\node[fix, right=2.5cm of pg] (pgfix) {Восстановление demo\\pg\_hba.conf\\listen\_addresses=localhost};

% Стрелки атаки (вниз по левой колонке)
\draw[arrow] (ssh) -- node[left, font=\small, align=center]{развитие\\доступа} (cacti);
\draw[arrow] (cacti) -- node[left, font=\small, align=center]{подключение\\к БД} (pg);

% Стрелки устранения (горизонтальные)
\draw[arrow, dashed, green!60!black] (ssh) -- node[above, font=\small]{устранение} (sshfix);
\draw[arrow, dashed, green!60!black] (cacti) -- node[above, font=\small]{устранение} (cactifix);
\draw[arrow, dashed, green!60!black] (pg) -- node[above, font=\small]{устранение} (pgfix);

\end{tikzpicture}
\caption{Схема атаки и последующего устранения последствий}
\label{fig:attack-fix}
\end{figure}

\clearpage

\section*{6. Реагирование и устранение последствий}
\addcontentsline{toc}{section}{6. Реагирование и устранение последствий}

\subsection*{6.1 Инцидент SSH Server}
\addcontentsline{toc}{subsection}{6.1 Инцидент SSH Server}

Работа с инцидентом SSH Server началась с проверки узла 10.10.1.31, который в сценарии атаки выступал исходной или промежуточной точкой доступа во внутреннюю инфраструктуру. На сервере были проверены состояние пользовательской сессии, доступность административных команд и признаки изменения системной конфигурации.

\subsubsection*{Уязвимость}

Уязвимость заключалась в недостаточно защищённом удалённом доступе по SSH. В ходе проверки было установлено, что при наличии подходящих учётных данных злоумышленник мог войти на сервер и выполнять действия от имени локального пользователя. Дополнительным фактором риска являлось отсутствие жёсткого ограничения источников подключения.

\subsubsection*{Последствие}

Последствием компрометации стало закрепление злоумышленника на узле через пользовательское задание cron. В задании был указан запуск файла \texttt{/home/user/evil.conf} каждую минуту. Такой способ опасен тем, что вредоносная активность может возобновляться автоматически после завершения отдельного процесса.

\subsubsection*{Устранение}

На первом этапе были выполнены проверка сервера и смена пароля пользователя. В терминале были просмотрены системное состояние и файл \texttt{/etc/hosts.deny}, проверена возможность использования \texttt{sudo} и выполнен переход к смене пароля. Поскольку \texttt{sudo} на данном узле отсутствовал, дальнейшие административные действия выполнялись через повышение прав штатными средствами системы.

\begin{figure}[H]
\centering
\includegraphics[width=0.9\textwidth]{\detokenize{screen_01.png}}
\caption{- Проверка SSH Server и смена пароля пользователя}
\end{figure}

На следующем этапе был проверен пользовательский crontab. В нём была обнаружена строка запуска \texttt{/home/user/evil.conf}. Это задание необходимо было удалить, а сам файл - изолировать или удалить после фиксации его имени и пути.

\begin{figure}[H]
\centering
\includegraphics[width=0.9\textwidth]{\detokenize{screen_02.png}}
\caption{- Обнаружение подозрительного задания cron на SSH Server}
\end{figure}

Основные действия по устранению последствий для SSH Server:
\begin{CodeBlock}
crontab -l
crontab -e
# удалить строку запуска /home/user/evil.conf
rm -f /home/user/evil.conf
passwd user
nano /etc/hosts.deny
systemctl status ssh
\end{CodeBlock}

После удаления задания cron и смены пароля был снижен риск повторного автоматического запуска вредоносного файла. Дополнительно необходимо проверить журналы \texttt{/var/log/auth.log} или их аналог в используемой системе, чтобы установить время входа и источник подключения.

В конце реагирования карточка инцидента по SSH Server была оформлена в системе Ampire.

\begin{figure}[H]
\centering
\includegraphics[width=0.9\textwidth]{\detokenize{screen_14.png}}
\caption{- Карточка инцидента SSH Server в Ampire}
\end{figure}

\subsection*{6.2 Инцидент Cacti}
\addcontentsline{toc}{subsection}{6.2 Инцидент Cacti}

Далее рассматривался инцидент Cacti на узле 10.10.2.32. В этой инфраструктуре Cacti выполнял роль системы мониторинга, поэтому имел доступ к данным о сетевых устройствах и мог запускать операции, связанные с их опросом. После компрометации SSH Server злоумышленник получил возможность развить атаку в сторону внутренних сервисов, включая Cacti.

\subsubsection*{Уязвимость}

Уязвимость была связана с наличием тестового device в Cacti. При разборе этого инцидента было отмечено, что такие объекты часто создаются для проверки настроек мониторинга, но после завершения настройки остаются в системе. Если к тестовому устройству привязаны нестандартные шаблоны, команды опроса или скрипты, оно может быть использовано как точка выполнения нежелательных действий.

\subsubsection*{Последствие}

Последствием эксплуатации стало получение административных прав на узле Cacti. В качестве индикатора компрометации было зафиксировано изменение файла sudoers, то есть изменение правил, определяющих, кто и какие команды может выполнять через \texttt{sudo}. Для сервера мониторинга это особенно критично, так как повышение привилегий позволяет злоумышленнику влиять не только на Cacti, но и на связанные системные службы.

\subsubsection*{Устранение}

Для устранения последствий на Cacti были выполнены следующие действия:
\begin{itemize}
  \item выявление и удаление тестового device, через который злоумышленник выполнил атаку;
  \item проверка списка устройств и шаблонов Cacti на наличие неиспользуемых объектов;
  \item восстановление файла \texttt{/etc/sudoers} и содержимого каталога \texttt{/etc/sudoers.d/};
  \item проверка прав пользователя, от имени которого работает Cacti;
  \item перезапуск веб-сервера и служб мониторинга после исправления конфигурации;
  \item фиксация результата в карточке инцидента Ampire.
\end{itemize}

Типовой набор команд администратора для проверки узла Cacti:
\begin{CodeBlock}
sudo visudo -c
sudo ls -la /etc/sudoers.d/
sudo grep -R "NOPASSWD\|cacti\|www-data" /etc/sudoers /etc/sudoers.d/
sudo systemctl restart apache2
sudo systemctl restart cron
\end{CodeBlock}

Ключевой профилактической мерой является удаление неиспользуемых devices и шаблонов из Cacti, а также запрет выполнения административных команд от имени веб-пользователя без строгой необходимости.

После выполнения работ инцидент был зарегистрирован в Ampire как атака на Data Center на базе Cacti. В карточке были указаны источник 10.10.1.31 (SSH Server), поражённый актив 10.10.2.32 (Cacti), описание эксплуатации тестового device и рекомендация удалить неиспользуемые устройства, уязвимые для атак.

\begin{figure}[H]
\centering
\includegraphics[width=0.9\textwidth]{\detokenize{screen_11.png}}
\caption{- Карточка инцидента Cacti в Ampire}
\end{figure}

\subsection*{6.3 Инцидент PostgreSQL}
\addcontentsline{toc}{subsection}{6.3 Инцидент PostgreSQL}

Последним был разобран инцидент PostgreSQL на сервере баз данных 10.10.2.26. Этот узел являлся критичным активом, так как хранил данные прикладной системы. В ходе расследования было установлено, что атака развивалась из внутренней сети, а в журналах PostgreSQL присутствовали признаки выполнения опасных SQL-команд и попыток запуска команд операционной системы.

\subsubsection*{Уязвимость}

Уязвимость была связана с небезопасной конфигурацией PostgreSQL. В файле \texttt{/etc/postgresql/13/main/pg\_hba.conf} было обнаружено широкое правило удалённого доступа, а сам сервер принимал подключения не только с локального узла. Это позволяло обращаться к PostgreSQL из внутренней сети и использовать скомпрометированные или слабые учётные данные.

При анализе журнала PostgreSQL была выявлена подозрительная активность: подключения из внутренней сети, обращения к системным объектам и попытка выполнить команду операционной системы через привилегированный механизм PostgreSQL. В журнале также присутствовало обращение к \texttt{pg\_shadow}, что указывает на попытку получить сведения об учётных записях базы данных.

\begin{figure}[H]
\centering
\includegraphics[width=0.9\textwidth]{\detokenize{screen_06.png}}
\caption{- Подозрительная активность в журнале PostgreSQL}
\end{figure}

\subsubsection*{Последствие}

Основным последствием атаки стало удаление базы данных \texttt{demo}. На сервере также был обнаружен подозрительный SQL-файл \texttt{hacker.sql}, расположенный в каталоге снимков PostgreSQL. Это указывает на то, что злоумышленник не только получил доступ к серверу баз данных, но и оставил артефакты, которые могли использоваться для повторного воздействия на систему.

\subsubsection*{Устранение}

Расследование началось с подключения к серверу PostgreSQL по SSH и проверки состояния баз данных. При попытке создать базу \texttt{demo} возникла ошибка, связанная с активными подключениями к шаблонной базе \texttt{template1}. Для устранения блокировки были завершены лишние backend-процессы PostgreSQL, после чего база \texttt{demo} была создана повторно и восстановлена из снимка.

\begin{figure}[H]
\centering
\includegraphics[width=0.9\textwidth]{\detokenize{screen_03.png}}
\caption{- Подключение к PostgreSQL и восстановление базы demo}
\end{figure}

Затем был выполнен поиск файлов резервных копий и SQL-снимков. В каталоге \texttt{/var/lib/postgresql/13/main/pg\_snapshots/} были найдены файлы \texttt{demo-snapshot.sql} и \texttt{hacker.sql}. Файл \texttt{demo-snapshot.sql} использовался как легитимный снимок для восстановления базы, а \texttt{hacker.sql} был отмечен как подозрительный артефакт, требующий изоляции и анализа.

\begin{figure}[H]
\centering
\includegraphics[width=0.9\textwidth]{\detokenize{screen_04.png}}
\caption{- Проверка баз PostgreSQL и поиск файлов резервных копий}
\end{figure}

\begin{figure}[H]
\centering
\includegraphics[width=0.9\textwidth]{\detokenize{screen_05.png}}
\caption{- Обнаружение SQL-снимков в каталоге pg\_snapshots}
\end{figure}

Для восстановления работоспособности была выполнена следующая последовательность:
\begin{CodeBlock}
ssh admin@10.10.2.26
sudo -u postgres psql
SELECT pg_terminate_backend(pid)
FROM pg_stat_activity
WHERE datname = 'template1'
  AND pid <> pg_backend_pid();
\q
sudo -u postgres createdb demo
sudo -u postgres psql -d demo -f /var/lib/postgresql/13/main/pg_snapshots/demo-snapshot.sql
\end{CodeBlock}

После восстановления базы были закрыты использованные уязвимости в конфигурации. В файле \texttt{/etc/postgresql/13/main/pg\_hba.conf} было закомментировано широкое правило удалённого доступа, а локальные подключения были ограничены Unix-сокетом и loopback-адресами с методом \texttt{md5}.

\begin{figure}[H]
\centering
\includegraphics[width=0.9\textwidth]{\detokenize{screen_08.png}}
\caption{- Исправление правил доступа в pg\_hba.conf}
\end{figure}

В файле \texttt{/etc/postgresql/13/main/postgresql.conf} параметр \texttt{listen\_addresses} был установлен в значение \texttt{localhost}. Это исключило приём сетевых подключений PostgreSQL со всей внутренней сети и оставило доступ только локальным процессам.

\begin{figure}[H]
\centering
\includegraphics[width=0.9\textwidth]{\detokenize{screen_09.png}}
\caption{- Ограничение адресов прослушивания PostgreSQL}
\end{figure}

После изменения конфигурации сервис PostgreSQL был перезапущен, а пароль пользователя \texttt{postgres} изменён. Это необходимо, так как журнал показывает попытки входа и выполнение действий от имени привилегированных пользователей базы данных.

\begin{figure}[H]
\centering
\includegraphics[width=0.9\textwidth]{\detokenize{screen_07.png}}
\caption{- Перезапуск PostgreSQL и смена пароля пользователя postgres}
\end{figure}

Итоговые команды усиления конфигурации:
\begin{CodeBlock}
sudo nano /etc/postgresql/13/main/pg_hba.conf
sudo nano /etc/postgresql/13/main/postgresql.conf
sudo systemctl restart postgresql
sudo -u postgres psql
ALTER USER postgres WITH PASSWORD '<новый_сложный_пароль>';
\q
\end{CodeBlock}

Таким образом, последствия атаки были устранены: база данных восстановлена из снимка, небезопасный сетевой доступ отключён, правила аутентификации исправлены, пароль суперпользователя обновлён, а подозрительная активность зафиксирована в карточке инцидента.

После завершения технических действий инцидент PostgreSQL был зарегистрирован в Ampire как атака на узел Data Center на базе Postgres. В карточке были указаны источник 10.10.1.31 (SSH Server), поражённый актив 10.10.2.26 (Postgres), описание последствий и рекомендация исправить файлы конфигурации.

\begin{figure}[H]
\centering
\includegraphics[width=0.9\textwidth]{\detokenize{screen_10.png}}
\caption{- Карточка инцидента PostgreSQL в Ampire}
\end{figure}

\clearpage

\section*{Заключение}
\addcontentsline{toc}{section}{Заключение}

В ходе научно-исследовательской практики был разобран комплексный инцидент информационной безопасности в учебной инфраструктуре. Атака затронула три ключевых компонента: SSH Server, Cacti и PostgreSQL. Для каждого компонента были определены признаки компрометации, возможный вектор атаки и меры устранения последствий.

На SSH Server был обнаружен механизм закрепления через cron и выполнена смена пароля пользователя. Для Cacti была зафиксирована эксплуатация тестового device с последующим изменением sudoers; в качестве меры защиты указаны удаление неиспользуемых объектов мониторинга и восстановление правил привилегий. На PostgreSQL были обнаружены следы удалённого подключения, обращения к чувствительным системным объектам и удаление базы данных. База \texttt{demo} была восстановлена из снимка, а конфигурация PostgreSQL усилена через ограничение \texttt{listen\_addresses} и исправление \texttt{pg\_hba.conf}.

По итогам практики описан порядок реагирования, который позволяет последовательно выполнять фиксацию инцидента, локализацию атаки, восстановление сервисов и профилактическое усиление инфраструктуры. Практическая ценность работы заключается в том, что действия администратора описаны с привязкой к конкретным узлам, адресам, файлам конфигурации и подтверждающим скриншотам.

\clearpage

\section*{Список литературы}
\addcontentsline{toc}{section}{Список литературы}

\begin{enumerate}
  \item PostgreSQL Documentation. Client Authentication. URL: \url{https://www.postgresql.org/docs/current/client-authentication.html}
  \item PostgreSQL Documentation. Server Configuration. URL: \url{https://www.postgresql.org/docs/current/runtime-config.html}
  \item PostgreSQL Documentation. SQL COPY. URL: \url{https://www.postgresql.org/docs/current/sql-copy.html}
  \item The OpenSSH Project. Manual pages. URL: \url{https://www.openssh.com/manual.html}
  \item Debian Administrator's Handbook. Service management and system administration. URL: \url{https://www.debian.org/doc/manuals/debian-handbook/}
  \item Cacti Documentation. Devices and templates. URL: \url{https://docs.cacti.net/}
  \item MITRE ATT\&CK. Persistence: Scheduled Task/Job. URL: \url{https://attack.mitre.org/techniques/T1053/}
  \item MITRE ATT\&CK. Valid Accounts. URL: \url{https://attack.mitre.org/techniques/T1078/}
  \item MITRE ATT\&CK. Account Discovery. URL: \url{https://attack.mitre.org/techniques/T1087/}
  \item ГОСТ Р 57580.1-2017. Безопасность финансовых операций. Защита информации финансовых организаций. Базовый состав организационных и технических мер.
\end{enumerate}

\end{document}
